\section{Интерактивный прототип}
\headingtextsep

Интерактивный прототип предназначен для проверки базовых сценариев взаимодействия пользователя с системой и навигации между основными экранами.

\textheadingsep
\subsection{Способ реализации}
\headingtextsep

Интерактивный прототип реализован в виде набора взаимосвязанных HTML-страниц с единым CSS-оформлением и минимальной логикой на JavaScript. Такой подход позволяет показать переходы между экранами, работу основных элементов навигации и наличие контекстных инструкций.

\textheadingsep
\subsection{Минимальный состав экранов}
\headingtextsep

Интерактивный прототип включает три взаимосвязанных экрана:

\begin{itemize}
\item главный экран проекта;
\item экран списка задач;
\item экран базы знаний.
\end{itemize}

Скриншоты этих экранов представлены в приложении Б на рисунках \ref{fig:appendix-int-overview}, \ref{fig:appendix-int-tasks} и \ref{fig:appendix-int-docs}.

\textheadingsep
\subsection{Навигация и интерактивность}
\headingtextsep

В интерактивном прототипе реализованы переходы между основными экранами через верхнее и боковое меню. Внутри экранов также представлены простые действия: выбор задачи из списка, фильтрация задач по состоянию и переключение статьи в разделе базы знаний. Эти действия не требуют полноценной бизнес-логики, но позволяют показать связное пользовательское взаимодействие.

\textheadingsep
\subsection{Контекстные инструкции}
\headingtextsep

На каждом экране предусмотрены краткие инструкции по работе с интерфейсом. На главном экране дано пояснение по структуре рабочего пространства, на экране задач --- подсказка по использованию фильтров и просмотру карточки задачи, на экране базы знаний --- инструкция по переходу между разделами и просмотру материалов. Наличие таких подсказок показано на скриншотах, включенных в приложение Б.